\chapter{Coordination and Messages}
\label{ch:coordination}

The two agents form a team, and all of their cooperation flows through the game's
chat channel as structured messages. This chapter describes the transport, the
message vocabulary, and the protocols built on top of it. The same machinery serves
two roles: the LLM coordinator uses it to push commands and policies to the executor,
and both agents use it to continuously share state and resolve conflicts.

\section{Agent Communication}
\label{sec:agent-comm}

Messages are exchanged over two primitives: \texttt{emitSay} delivers a message
privately to a single recipient, while \texttt{emitShout} broadcasts to everyone. The
team prefers the private channel, a directed \texttt{emitSay} to the partner, and
falls back to a shout only if the direct send fails. This keeps the public chat
uncluttered and avoids leaking coordination to competing teams. Incoming chat is
pre-filtered: each agent ignores its own messages and anything that does not match
the cooperative JSON schema, so ordinary Admin text and opponent chatter never reach
the protocol handlers.

Underlying everything is a steady stream of \texttt{PEER\_STATUS} heartbeats sent
directly to the partner a few times per second. Each carries the sender's position,
score, carried inventory, planned path and next step, and known crates. This keeps each
agent's model of its partner current even when it is out of sensing range, shares crate
sightings so both navigate around the same obstacles, and lets each agent see in advance
when they are converging on the same tile, so one can yield or re-deliberate toward a
different target without any explicit negotiation or locking.

\section{Message Structure}
\label{sub:msg-structure}
Every cooperative message is a JSON object with a \texttt{type} field that selects a
handler; the remaining fields are the payload. The vocabulary groups into three
families:

\begin{itemize}
  \item \textbf{Status:} \texttt{PEER\_STATUS}, for position, cargo and obstacle
        sharing, and the \texttt{SYNC\_REQ}/\texttt{SYNC\_ACK} pair used once at
        startup so neither agent acts before its partner is online (\cref{sec:init_state}).
  \item \textbf{Direct commands} (coordinator to executor): \texttt{MOVE\_TO},
        \texttt{HOLD}, \texttt{RESUME}, \texttt{APPLY\_RULES}, \texttt{SET\_VARIABLE},
        and \texttt{INSTRUCT\_SAY}. Of these, only \texttt{MOVE\_TO} expects a reply:
        the executor returns a \texttt{MOVE\_TO\_ACK} once it has reached the target
        or failed to, so the coordinator knows the move is complete.
        \texttt{INSTRUCT\_SAY} is the coordinator asking the executor to relay the
        LLM's chat answer to the Admin, a redundant delivery path that makes the
        reply more likely to reach the Admin.
  \item \textbf{Contracts:} \texttt{PROPOSE\_CONTRACT}, \texttt{ACCEPT\_CONTRACT},
        \texttt{CLOSE\_CONTRACT} for the propose/accept/close lifecycle, plus
        \texttt{SIGNAL\_READY} and \texttt{RELEASE\_CARGO} as the status transitions
        inside the hand-off and relay choreography: \texttt{SIGNAL\_READY} announces
        that an agent has dropped its cargo and stepped aside so the partner may
        collect it, and \texttt{RELEASE\_CARGO} marks the cargo as released for the
        partner to take.
\end{itemize}

On receipt, a handler applies the message directly to the agent's own belief base.
This is the mechanism by which a coordinator command takes effect remotely: a
\texttt{MOVE\_TO} becomes a high-priority \texttt{admin\_move} contract (and
implicitly releases any standing hold), an \texttt{APPLY\_RULES} updates the policy
set, \texttt{HOLD}/\texttt{RESUME} toggle the paused state, and a \texttt{SET\_VARIABLE}
writes shared memory. The executor thus needs no special command interpreter, since
coordination and perception update the same beliefs through the same revision path.

\section{Coordination Workflow}
\label{sec:coord-workflow}
\paragraph{Contract handshake.} Multi-agent manoeuvres are negotiated as
\emph{contracts} rather than hard-coded. One agent proposes, the other validates that 
the target tile is valid and walkable and accepts, both mark it active, and either may cancel it
later. We chose a propose/accept handshake over simply commanding the partner so that
a contract is never acted upon unless both agents agree it is feasible, which avoids
one agent waiting forever for a partner that cannot reach the meeting point. Three
cooperation patterns are built on this handshake.

\begin{itemize}
  \item \textbf{Meet-and-wait} (\texttt{RENDEZVOUS}, and \texttt{CLEARING} as a
        special case): both agents converge within a radius of a tile and pause. We
        decided to model ``clear this gate'' as a rendezvous rather than a separate
        mechanism, because mechanically it is the same act, bring an agent to a point
        and hold, and collapsing the two keeps the executor simple.
  \item \textbf{Hand-off} (\texttt{HANDOFF}): one agent carries its cargo to a shared
        tile, drops it, and steps aside so the other can collect it. This exists for
        the case where the carrier cannot reach the delivery zone but its partner can.
  \item \textbf{Relay} (\texttt{RELAY}): a persistent contract in which a designated
        courier repeatedly farms parcels and drops them, while the partner delivers
        them. The decision driving this pattern is strategic: the game awards a bonus
        when a parcel is picked up by one agent and delivered by another, and the
        relay is how the team deliberately manufactures that situation. A consequence
        we had to handle explicitly is the choice of drop tile. It is placed
        \emph{adjacent to} the best delivery zone but never \emph{on} it, because a
        courier putting a parcel down on a delivery tile would score the delivery
        itself and so cancel the very bonus the relay is meant to earn.
\end{itemize}

Blocking commands wait for confirmation before the coordinator proceeds: a directed
move resolves only once a \texttt{MOVE\_TO\_ACK} or a position match is observed, a
deliberate choice so the LLM's reasoning never advances on an unconfirmed physical
effect. Control actions such as \texttt{RESUME} clear outstanding contracts so no
agent is left waiting in a coordination loop.
